\section{Quadratic perspectives}

\begin{lemma}[Finite quadratic-perspective inequality]
    \label{lem:finite_quadratic_perspective}
    \lean{Matrix.sum_conjTranspose_inv_mul_le,
        Matrix.conjTranspose_inv_mul_le_of_concave}
    \leanok
    Let $R_i>0$ and $S>0$ be square matrices, let $L_i$ be rectangular
    matrices of compatible sizes, and let $w_i\geq0$ be finitely many weights.
    If $\sum_i w_iR_i\leq S$, then
    \[
        \left(\sum_i w_iL_i\right)^\dagger S^{-1}
        \left(\sum_i w_iL_i\right)
        \leq \sum_i w_iL_i^\dagger R_i^{-1}L_i.
    \]
    There is no normalization condition on the weights.
\end{lemma}

\begin{proof}\leanok
    The Schur complement shows that each matrix
    \[
        \begin{pmatrix}L_i^\dagger R_i^{-1}L_i&L_i^\dagger\\L_i&R_i\end{pmatrix}
    \]
    is positive semidefinite. Take their weighted sum and add the positive
    block diagonal matrix with diagonal blocks $0$ and $S-\sum_iw_iR_i$.
    The lower-right block is now $S$. Its Schur complement gives the result.
\end{proof}

\begin{theorem}[Quadratic perspective with a concave denominator]
    \label{thm:quadratic_perspective_concave_denominator}
    \lean{Matrix.convexOn_quadraticPerspective, Matrix.convexOn_affine_quadraticPerspective}
    \leanok
    Let $D$ be a convex subset of a real vector space, and let $g$ be
    operator concave on $D$ with $g(R)>0$ for every $R\in D$.
    Then $(L,R)\mapsto L^\dagger g(R)^{-1}L$ is jointly operator convex
    on the product of the rectangular matrix space and $D$.
    More generally, if $h$ is an affine rectangular-matrix-valued map on $D$,
    then $R\mapsto h(R)^\dagger g(R)^{-1}h(R)$ is operator convex.
    For Hermitian $L$ this is the quadratic case of the perspective in
    \cite[Theorem~5.14]{Wolf2012Quantum}. The assertion here concerns this
    quadratic case; the theorem for a general operator convex function is
    a separate result.
\end{theorem}

\begin{proof}\leanok
    \uses{lem:finite_quadratic_perspective}
    For $a,b\geq0$ with $a+b=1$, concavity gives
    $ag(R_1)+bg(R_2)\leq g(aR_1+bR_2)$. Apply
    Lemma~\ref{lem:finite_quadratic_perspective} with these two weights.
    An affine numerator preserves the convex combination used in this argument.
    For Hermitian $L$, cancellation of the square roots in
    $g(R)^{1/2}(g(R)^{-1/2}Lg(R)^{-1/2})^2g(R)^{1/2}$ gives
    $Lg(R)^{-1}L$.
\end{proof}

\begin{corollary}[Inverse-power quadratic perspectives]
    \label{cor:inverse_power_quadratic_perspective}
    \lean{Matrix.conjTranspose_rpow_neg_mul_convex,
        Matrix.convexOn_conjTranspose_rpow_neg_mul,
        Matrix.sum_conjTranspose_rpow_neg_mul_le}
    \leanok
    For $0\leq\alpha\leq1$, the map
    $(L,R)\mapsto L^\dagger R^{-\alpha}L$ is jointly operator convex
    for arbitrary rectangular $L$ and positive definite $R$.
    In particular, for finite probability weights $w_i$,
    \[
        \left(\sum_iw_iL_i\right)^\dagger
        \left(\sum_iw_iR_i\right)^{-\alpha}
        \left(\sum_iw_iL_i\right)
        \leq\sum_iw_iL_i^\dagger R_i^{-\alpha}L_i.
    \]
    For positive $L$ this is the corrected-domain form of
    \cite[Corollary~5.4(1)]{Wolf2012Quantum}. The printed Hermitian-denominator
    condition is insufficient; positivity and invertibility restore the
    preceding perspective theorem
    \cite{gap:wolf_ch5_quadratic_perspective_domain}.
    Singular denominators with negative powers are not included.
\end{corollary}

\begin{proof}\leanok
    \uses{thm:quadratic_perspective_concave_denominator}
    The power $R\mapsto R^\alpha$ is operator concave on positive matrices
    for $0\leq\alpha\leq1$, and is positive definite when $R$ is.
    Moreover $(R^\alpha)^{-1}=R^{-\alpha}$. Apply
    Theorem~\ref{thm:quadratic_perspective_concave_denominator}.
    Finite convex combinations give the displayed inequality.
\end{proof}

\section{Two-factor tensor powers}

\begin{lemma}[Tensor quadratic forms and traces]
    \label{lem:vec_transpose_kronecker_quadratic}
    \lean{Matrix.vec_transpose_kronecker_quadratic}
    \leanok
    For complex square matrices $A,B,K$ of the same size, put
    $v_K=\sum_{i,j}K_{ij}\lvert i\rangle\otimes\lvert j\rangle$. Then
    \[
        \langle v_K,(A\otimes B)v_K\rangle
        =\tr(K^\dagger A K B^T).
    \]
    This is \cite[Equation~(5.48)]{Wolf2012Quantum}. No positivity assumption
    is needed for the identity.
\end{lemma}

\begin{proof}\leanok
    Expand the quadratic form in the product basis. Its summands are
    $\overline{K_{ij}}A_{ik}B_{jl}K_{kl}$, which also give the trace on the right.
\end{proof}

\begin{theorem}[Joint concavity of two tensor powers]
    \label{thm:rpow_kronecker_rpow_concave}
    \lean{Matrix.rpow_kronecker_rpow_concave}
    \leanok
    For $x,y\geq0$ with $x+y\leq1$, the map
    $(A,B)\mapsto A^x\otimes B^y$ is jointly operator concave on pairs of
    positive semidefinite matrices of the same size. This is
    \cite[Corollary~5.5(1)]{Wolf2012Quantum}.
\end{theorem}

\begin{proof}\leanok
    \uses{lem:vec_transpose_kronecker_quadratic, thm:lieb_concavity_subboundary}
    Every vector of the tensor product is $v_K$ for some $K$.
    The preceding identity identifies its quadratic form with
    $\tr(K^\dagger A^xK(B^T)^y)$, since real powers commute with transpose.
    The concave Ando--Lieb inequality proves the desired inequality for every
    quadratic form, hence in the operator order.
\end{proof}

\begin{theorem}[Joint convexity of two tensor powers]
    \label{thm:rpow_kronecker_rpow_convex}
    \lean{Matrix.rpow_kronecker_rpow_convex}
    \leanok
    For $1\leq x\leq2$ and $1-x\leq y\leq0$, the map
    $(A,B)\mapsto A^x\otimes B^y$ is jointly operator convex for $A\geq0$
    and $B>0$, of the same size. Taking $x=1+\alpha$ and
    $y=-\alpha\beta$ gives \cite[Corollary~5.5(3)]{Wolf2012Quantum}
    for $\alpha,\beta\in[0,1]$, on this invertible domain.
\end{theorem}

\begin{proof}\leanok
    \uses{lem:vec_transpose_kronecker_quadratic, thm:lieb_convexity_posDef}
    Apply the convex Ando--Lieb inequality with the second input transposed.
    The vectorization identity tests the resulting operator inequality on
    every vector. Transpose preserves positive definiteness, so all inverse
    powers are defined on the required domain.
\end{proof}
